% 12-institutions.tex: ткань обмена.
\section{Учреждения: обмен, время и подписи без журнала сообщений}
\label{sec:institutions}

\subsection{Выворачивание наизнанку}

Обыкновенная ткань обмена между учреждениями делает обмен доказуемым, записывая обмениваемые
сообщения: посредующий узел подписывает трафик и помечает его временем, чтобы расписку можно было
позже показать третьей стороне. Свидетельство при этом настоящее, и цена его тоже настоящая:
хранилище того, кто о ком что спрашивал, а это ровно та запись в масштабе населения, которую
призвание запрещает. Ткань обмена Polaris есть выворачивание этого наизнанку. Она доказывает те же
три вещи, что доказывает журнал сообщений, а именно что обмен состоялся, между этими сторонами и с
дозволения, не удерживая при этом ни байта полезной нагрузки, поскольку фиксирует запрос и ответ
хешем и никогда содержимым. Всякий примитив в ней, подписи, отделяемый верификатор, журналы
прозрачности, граф доверия, уже существовал; ткань их складывает. \sImpl

\subsection{Реестр: обнаружение, а не настройка}

Потребитель узнаёт, что экземпляр предлагает и чему доверяет, из одного подписанного предмета:
форматы и версии протокола, на которых он говорит, и принимаемые им алгоритмы, род каждой службы,
шаблон её пути и способ опознания, журналы прозрачности, каждый федеративный орган с его реестром
ключей, контексты и доказательство, требуемое каждым из них, граф доверия внутри контекста и
доверяющие стороны, но только по имени и области. Всякий факт есть представление над слоем органов;
реестр не добавляет собственной истины. Он обязан быть подписан тем действующим ключом, который сам
же перечисляет для своего издателя, так что самозванец не способен опубликовать реестр от имени
органа, и он нетранзитивен: он описывает экземпляр своего издателя и никогда чужой. Учение на двух
экземплярах читает путь службы из забранного реестра и вызывает её. Ни токена, ни держателя, ни
записи о проверке там не появляется: карта учреждений, но никогда карта людей. Одно, чего реестр не
несёт, записано в разделе~\ref{sec:signatures}: принимаемые алгоритмы появляются голым перечнем
имён без состояния, так что читающий его верификатор не способен узнать, что какой-то из них
объявлен устаревшим. \sImpl

\subsection{Шлюз и расписка}

Расписка фиксирует ключ запрашивающего, отвечающего, контекст, два хеша тел, ту аттестацию,
которая дозволила запрашивающему, и момент. Конечная точка чеканки принимает только хеши, так что
приложение не способно удержать того, чего ему никогда не дают. Третья сторона проверяет автономно,
что расписка подлинна и что запрашивающему было дозволено в этом контексте, исходя из описей,
которым она доверяет. Сторона, держащая тело, способна подтвердить, что хеш его связывает; сторона,
которая его не держит, всё равно держит подписанный отчёт отвечающего. Одна поправка из внешнего
ревью принадлежит этому месту: одна расписка не способна доказать, что запрашивающий участвовал,
потому что нечестный отвечающий может подписать любую расписку, какая ему нравится; собственный
подписанный конверт запрашивающего рядом с ней способен, и спецификация формата передачи называет
свидетельством именно пару. Множество расписок само есть журнал прозрачности, так что учреждение не
способно позже отрицать, что расписка существовала, и всякий способен увидеть, сколько их было
вычеканено, тогда как каждая расписка остаётся в руках своих сторон. \sImpl

\figfile{gateway}{Ведомственный обмен через шлюз. Запрашивающий подписывает конверт своим
зарегистрированным ключом; собственный экземпляр отвечающего опознаёт его по ключу, который уже
знает, дозволяет его через собственную аттестацию в контексте конверта ещё до того, как что-либо
покинет процесс, поглощает одноразовое число в регистре повторов, работающем только на добавление,
передаёт тело лишь той вышестоящей службе, которую настроил оператор, и возвращает ответ с
подписанной распиской. Не пишется ничего, кроме хеша расписки и поглощённого одноразового числа.
Свидетельством обмена служит пара из конверта и расписки, проверяемая автономно третьей стороной и
без всякого тела.}{fig:gateway}

Дозволение на шлюзе направленно. Начиная с \code{v9.333} само отвечающее ведомство обязано
аттестовать ключ запрашивающего в контексте конверта; аттестация любым другим ведомством на этом
экземпляре не дозволяет там ничего. До этого исправления шлюз принимал аттестацию от кого угодно на
экземпляре, а это было то самое транзитивное доверие, которое слой федерации отвергает, вновь
всплывшее уровнем выше. \sImpl

\subsection{Время, которое ничего не узнаёт}

Служба отметок времени привязывает произвольный дайджест SHA3-256 к моменту под зарегистрированным
ключом органа. Она принимает дайджест, но никогда содержимое, и не держит записи по каждому
запросу, потому что служба отметок времени, записывающая каждый запрос, есть хранилище того, кто
что и когда пометил временем. Свидетельство есть подписанное утверждение в руках запрашивающего.
Начиная с \code{v9.341} вызывающий вправе выбрать заякоривание: хеш самой отметки времени входит в
засвидетельствованный журнал только на добавление, а свидетельство включения возвращается
подшитым, так что отметка времени, проставленная задним числом под украденным ключом органа, есть
отметка, отсутствующая во всякой засвидетельствованной вершине той эпохи, на которую она
притязает. Орган при этом удерживает один дайджест и один момент, то есть запись об объёме и
времени, но никогда о содержимом, а запрос по умолчанию по-прежнему не оставляет ничего. \sImpl

\subsection{Подписать документ так, чтобы он пережил ключ}

Держатель в Polaris вправе владеть ключом (раздел~\ref{sec:agents}), и вправе не владеть, поэтому
модель подписания документов остаётся нотариальной: выдающий орган подписывает от имени держателя,
только что доказавшего владение, ровно так же, как утверждает для него состояние, и записывает,
кто это был, по хешу. Подпись органа говорит, что держатель этого удостоверения, действующего на
этот момент, попросил подписать этот дайджест, а вложенные свидетельства позволяют всякому
проверить это утверждение позже. \sImpl

Долгосрочная проверка есть место, где внешнее ревью и собственные учения репозитория заострили
утверждение. Подлинная отметка времени не есть доверенная, поскольку подписать её может всякий, а
собственная отметка времени подписанта проставима задним числом тем, кто держит ключ. Поэтому
сильное утверждение требует, чтобы отметка времени была доверена верификатором, поставлена ключом,
отличным от ключа подписанта, действующим на тот момент по списку доверия, и, если верификатор того
просит, заякорена в засвидетельствованном журнале либо подтверждена кворумом независимых органов.
Верификатор, которому не дали якорей отметки времени, сообщает факты и не утверждает ничего.
Проверка якорей есть и в отделяемом верификаторе, и в обоих SDK, так что реализация, не
импортирующая ни строки кода Polaris, способна решить вопрос о долгосрочной действительности.
\sImpl

\figfile{signing-chain}{Подписанный документ и его долгосрочные свидетельства. Документ никогда не
покидает своего держателя; его дайджест подписывает учреждение либо выдающий орган держателя от
имени держателя после доказательства владения, с записью по хешу удостоверения. Свидетельства,
закреплённые в момент подписания, а именно отметка времени от отличного ключа, опись подписанта,
его контрольная точка и лента, вместе с собственным списком доверия верификатора позволяют
верификатору решить вопрос о действительности на тот момент спустя долгое время после вывода
подписывающего ключа из обращения. Подпись, момент которой приходится на время после
компрометации, отвергается, даже если вложенная в неё опись говорит, что ключ был действующим.}{fig:signing-chain}

\begin{layercard}{Карточка слоя: учреждения}
\qa{Что это}{Подписанный реестр, шлюз обмена, расписка обмена и её журнал, служба отметок времени с необязательным заякориванием, подписание документов с долгосрочной проверкой, список доверия.}
\qa{Зачем оно существует}{Настоящим программным экосистемам нужны обнаружение, посредованные запросы, подписи и время. Каждое из этого принято строить на журнале содержимого, а призвание такой журнал запрещает.}
\qa{Чему доверяет}{Зарегистрированным ведомственным ключам; графу доверия внутри контекста; собственным якорям отметок времени и списку доверия верификатора.}
\qa{Чему отказывает}{Телу на конечной точке чеканки; адресу, названному запросом; переигранному одноразовому числу; аттестации, сделанной не самим отвечающим; отметке времени собственным ключом подписанта как сильному свидетельству; подписи-заглушке как способу опознания.}
\qa{Что видит}{Дайджесты, ключи, контексты, моменты.}
\qa{Чего не видит}{Обмениваемых сообщений, помеченного временем содержимого, подписанного документа: ничего из этого ей не дают никогда.}
\qa{Если откажет}{Поддельная расписка падает на подписи отвечающего; расписке нечестного отвечающего противоречит отсутствующий конверт; украденный ключ отметок времени ловится засвидетельствованным якорем или кворумом; скомпрометированный ключ подписанта обессмысливает лишь подписанное им после того момента.}
\qa{Кем проверяется}{\code{check\_registry}, \code{check\_exchange\_gateway}, \code{check\_exchange\_receipt}, \code{check\_timestamp\_authority}, \code{check\_document\_signing}, \code{check\_ltv\_timestamp\_trust}, \code{check\_broker\_policy\_bound}.}
\qa{Сейчас или потом}{Сейчас, между двумя экземплярами. Настоящие учреждения потом.}
\qa{Внутри и снаружи}{Внутри: прозрачность, наблюдающая за журналами расписок и отметок времени. Снаружи: эксплуатация, которая держит всё это на ходу.}
\end{layercard}

\nextlayer{Всякий слой до сих пор есть утверждение о правильности. Всё это не имеет значения, если
система валится под нагрузкой, теряет базу данных или не может быть развёрнута никем, кроме
автора. Следующий слой есть выживание.}
